% Generated by tools/paper/render.py from supplement.py; edit the source there, not this file.
\documentclass[10pt]{article}
\usepackage[a4paper,margin=1.8cm,columnsep=0.7cm]{geometry}
\usepackage[T1]{fontenc}
\usepackage{lmodern}
\usepackage{microtype}
\usepackage{amsmath,amssymb}
\usepackage{booktabs}
\usepackage{array}
\usepackage{longtable}
\usepackage{graphicx}
\usepackage{caption}
\captionsetup{font=small,labelfont=bf}
\usepackage{xcolor}
\usepackage{tikz}
\usepackage[round]{natbib}
\usepackage[hidelinks]{hyperref}
\usepackage{url}

\renewcommand{\thesection}{S\arabic{section}}
\renewcommand{\thetable}{S\arabic{table}}
\begin{document}
{\Large\bfseries Supplementary material for pgphase: pangenome graph-aware phasing of long reads with known and sample-specific variants\par}\medskip
{\large Konstantinos Kyriakidis and Benedict Paten}\par\bigskip
\section{Per-chromosome results}

Table S1 gives, per chromosome and tool, the phaseable reads on the arms by class, the discordant reads whose parental assignment is confident, the switch and flip errors of a record-level comparison and the phased truth heterozygotes. The MHC lies inside chr6.

{\footnotesize
\setlength{\LTcapwidth}{\linewidth}
\begin{longtable}{@{}lrrrrrrr@{}}
\caption{Phaseable reads on the arms (correct, discordant, unphased), discordant reads whose parental assignment is confident (HapQ $\ge$ 30), arm switches and flips, and phased truth heterozygotes, per chromosome and tool.}\label{tab:perchrom}\\
\toprule
Tool & Correct & Discordant & Unphased & Disc. HapQ$\ge$30 & Switches & Flips & Phased hets \\
\midrule
\endfirsthead
\caption[]{(continued)}\\
\toprule
Tool & Correct & Discordant & Unphased & Disc. HapQ$\ge$30 & Switches & Flips & Phased hets \\
\midrule
\endhead
\multicolumn{8}{@{}l}{\emph{chr12}} \\
pgphase & \textbf{523{,}469} & \textbf{2{,}392} & \textbf{356} & \textbf{380} & \textbf{4} & 14 & 134{,}155 \\
HiPhase & 503{,}601 & 17{,}847 & 4{,}769 & 14{,}458 & 23 & 19 & 127{,}723 \\
LongPhase & 493{,}458 & 4{,}086 & 28{,}673 & 2{,}612 & 10 & \textbf{12} & 109{,}900 \\
WhatsHap & 486{,}797 & 23{,}995 & 15{,}425 & 19{,}876 & 55 & 196 & 128{,}022 \\
WhatsHap-dg & 497{,}143 & 13{,}178 & 15{,}896 & 10{,}940 & 25 & 39 & 127{,}238 \\
longcallD & 485{,}719 & 5{,}534 & 34{,}964 & 3{,}161 & 26 & 27 & \textbf{136{,}655} \\
\multicolumn{8}{@{}l}{\emph{chr17}} \\
pgphase & \textbf{286{,}052} & \textbf{3{,}025} & \textbf{871} & \textbf{1{,}073} & \textbf{4} & 11 & 78{,}209 \\
HiPhase & 272{,}356 & 5{,}632 & 11{,}960 & 3{,}904 & 21 & 13 & 73{,}299 \\
LongPhase & 262{,}901 & 3{,}526 & 23{,}521 & 2{,}493 & 21 & \textbf{9} & 63{,}498 \\
WhatsHap & 254{,}920 & 13{,}660 & 21{,}368 & 11{,}175 & 71 & 127 & 73{,}219 \\
WhatsHap-dg & 258{,}127 & 10{,}383 & 21{,}438 & 8{,}465 & 42 & 30 & 72{,}338 \\
longcallD & 252{,}998 & 4{,}118 & 32{,}832 & 2{,}009 & 30 & 26 & \textbf{78{,}340} \\
\multicolumn{8}{@{}l}{\emph{chr18}} \\
pgphase & \textbf{295{,}091} & \textbf{1{,}513} & \textbf{198} & \textbf{322} & \textbf{9} & \textbf{7} & 73{,}464 \\
HiPhase & 281{,}144 & 11{,}124 & 4{,}534 & 9{,}247 & 23 & 11 & 71{,}236 \\
LongPhase & 273{,}676 & 2{,}367 & 20{,}759 & 1{,}493 & 14 & \textbf{7} & 61{,}960 \\
WhatsHap & 270{,}121 & 15{,}500 & 11{,}181 & 13{,}504 & 31 & 83 & 71{,}527 \\
WhatsHap-dg & 271{,}585 & 13{,}984 & 11{,}233 & 12{,}275 & 22 & 21 & 70{,}949 \\
longcallD & 272{,}104 & 3{,}026 & 21{,}672 & 1{,}740 & 27 & 11 & \textbf{75{,}660} \\
\multicolumn{8}{@{}l}{\emph{chr19}} \\
pgphase & \textbf{189{,}223} & \textbf{1{,}576} & \textbf{1{,}507} & \textbf{418} & \textbf{13} & \textbf{7} & 69{,}661 \\
HiPhase & 177{,}963 & 7{,}966 & 6{,}377 & 6{,}669 & 28 & 15 & 66{,}225 \\
LongPhase & 174{,}040 & 4{,}477 & 13{,}789 & 3{,}951 & 19 & 15 & 56{,}005 \\
WhatsHap & 164{,}506 & 17{,}662 & 10{,}138 & 15{,}884 & 49 & 104 & 66{,}329 \\
WhatsHap-dg & 172{,}415 & 9{,}715 & 10{,}176 & 8{,}566 & 30 & 29 & 65{,}729 \\
longcallD & 172{,}649 & 2{,}497 & 17{,}160 & 1{,}260 & 19 & 26 & \textbf{71{,}026} \\
\multicolumn{8}{@{}l}{\emph{chr20}} \\
pgphase & \textbf{217{,}593} & 1{,}250 & \textbf{443} & \textbf{174} & \textbf{2} & 10 & 64{,}312 \\
HiPhase & 209{,}183 & 6{,}240 & 3{,}863 & 5{,}088 & 16 & 12 & 62{,}431 \\
LongPhase & 203{,}827 & \textbf{1{,}018} & 14{,}441 & 318 & 8 & \textbf{7} & 54{,}894 \\
WhatsHap & 199{,}331 & 11{,}450 & 8{,}505 & 9{,}734 & 33 & 84 & 62{,}704 \\
WhatsHap-dg & 204{,}156 & 6{,}788 & 8{,}342 & 5{,}453 & 17 & 24 & 62{,}003 \\
longcallD & 200{,}294 & 2{,}478 & 16{,}514 & 1{,}225 & 15 & 14 & \textbf{66{,}157} \\
\multicolumn{8}{@{}l}{\emph{chr6}} \\
pgphase & \textbf{672{,}781} & \textbf{4{,}070} & \textbf{1{,}111} & \textbf{635} & \textbf{9} & 18 & 176{,}507 \\
HiPhase & 640{,}944 & 23{,}298 & 13{,}720 & 18{,}852 & 36 & 19 & 168{,}982 \\
LongPhase & 621{,}588 & 9{,}369 & 47{,}005 & 7{,}434 & 14 & \textbf{14} & 146{,}002 \\
WhatsHap & 621{,}507 & 26{,}366 & 30{,}089 & 21{,}908 & 61 & 248 & 169{,}239 \\
WhatsHap-dg & 625{,}581 & 21{,}691 & 30{,}690 & 18{,}226 & 40 & 52 & 168{,}322 \\
longcallD & 616{,}432 & 9{,}286 & 52{,}244 & 5{,}999 & 43 & 34 & \textbf{176{,}859} \\
\multicolumn{8}{@{}l}{\emph{chr16}} \\
pgphase & \textbf{287{,}337} & \textbf{3{,}068} & \textbf{1{,}641} & \textbf{1{,}295} & \textbf{9} & 11 & 83{,}256 \\
HiPhase & 266{,}620 & 7{,}305 & 18{,}121 & 5{,}628 & 21 & \textbf{7} & 79{,}853 \\
LongPhase & 267{,}580 & 3{,}955 & 20{,}511 & 2{,}988 & 10 & 8 & 70{,}833 \\
WhatsHap & 250{,}720 & 13{,}701 & 27{,}625 & 11{,}560 & 46 & 108 & 79{,}405 \\
WhatsHap-dg & 258{,}392 & 5{,}900 & 27{,}754 & 4{,}397 & 21 & 23 & 78{,}521 \\
longcallD & 250{,}511 & 3{,}675 & 37{,}860 & 1{,}994 & 27 & 27 & \textbf{84{,}142} \\
\multicolumn{8}{@{}l}{\emph{chr21}} \\
pgphase & \textbf{127{,}508} & 782 & \textbf{305} & \textbf{72} & \textbf{8} & \textbf{3} & 36{,}482 \\
HiPhase & 124{,}767 & 1{,}637 & 2{,}191 & 876 & 10 & 9 & 34{,}923 \\
LongPhase & 119{,}153 & \textbf{641} & 8{,}801 & 216 & 9 & 8 & 30{,}589 \\
WhatsHap & 120{,}306 & 3{,}006 & 5{,}283 & 1{,}897 & 18 & 47 & 34{,}870 \\
WhatsHap-dg & 120{,}924 & 2{,}213 & 5{,}458 & 1{,}491 & 11 & 10 & 34{,}577 \\
longcallD & 117{,}472 & 1{,}159 & 9{,}964 & 530 & 13 & 7 & \textbf{37{,}207} \\
\multicolumn{8}{@{}l}{\emph{MHC (chr6:28.5--34.0~Mb)}} \\
pgphase & \textbf{20{,}860} & 77 & \textbf{67} & 14 & \textbf{0} & \textbf{1} & \textbf{14{,}934} \\
HiPhase & 20{,}363 & 55 & 586 & \textbf{11} & 2 & \textbf{1} & 14{,}488 \\
LongPhase & 19{,}703 & 51 & 1{,}250 & 18 & \textbf{0} & \textbf{1} & 12{,}475 \\
WhatsHap & 19{,}937 & 108 & 959 & 66 & \textbf{0} & 16 & 14{,}508 \\
WhatsHap-dg & 20{,}002 & \textbf{47} & 955 & 16 & \textbf{0} & 3 & 14{,}449 \\
longcallD & 19{,}035 & 125 & 1{,}844 & 58 & 2 & \textbf{1} & 13{,}196 \\
\bottomrule
\end{longtable}
}

\section{Switch errors by haplotype sequence}

Table S2 classes every arm switch of Table S1 by the haplotype sequences around it (main text Section 2.8).

{\footnotesize
\setlength{\LTcapwidth}{\linewidth}
\begin{longtable}{@{}lrrrr@{}}
\caption{Arm switches per chromosome and tool, classified by the haplotype sequences around them (Section 2.8).}\label{tab:census}\\
\toprule
Tool & Switches & Real & Representation & Undecided \\
\midrule
\endfirsthead
\caption[]{(continued)}\\
\toprule
Tool & Switches & Real & Representation & Undecided \\
\midrule
\endhead
\multicolumn{5}{@{}l}{\emph{chr12}} \\
pgphase & 4 & 0 & 2 & 2 \\
HiPhase & 23 & 18 & 0 & 5 \\
LongPhase & 10 & 2 & 0 & 8 \\
WhatsHap & 55 & 28 & 4 & 23 \\
WhatsHap-dg & 25 & 16 & 0 & 9 \\
longcallD & 26 & 20 & 3 & 3 \\
\multicolumn{5}{@{}l}{\emph{chr17}} \\
pgphase & 4 & 0 & 2 & 2 \\
HiPhase & 21 & 13 & 0 & 8 \\
LongPhase & 21 & 10 & 1 & 10 \\
WhatsHap & 71 & 43 & 8 & 20 \\
WhatsHap-dg & 42 & 27 & 3 & 12 \\
longcallD & 30 & 16 & 6 & 8 \\
\multicolumn{5}{@{}l}{\emph{chr18}} \\
pgphase & 9 & 0 & 4 & 5 \\
HiPhase & 23 & 13 & 0 & 10 \\
LongPhase & 14 & 2 & 1 & 11 \\
WhatsHap & 31 & 21 & 1 & 9 \\
WhatsHap-dg & 22 & 13 & 0 & 9 \\
longcallD & 27 & 15 & 8 & 4 \\
\multicolumn{5}{@{}l}{\emph{chr19}} \\
pgphase & 13 & 7 & 1 & 5 \\
HiPhase & 28 & 22 & 0 & 6 \\
LongPhase & 19 & 9 & 3 & 7 \\
WhatsHap & 49 & 34 & 2 & 13 \\
WhatsHap-dg & 30 & 22 & 0 & 8 \\
longcallD & 19 & 14 & 0 & 5 \\
\multicolumn{5}{@{}l}{\emph{chr20}} \\
pgphase & 2 & 0 & 2 & 0 \\
HiPhase & 16 & 13 & 0 & 3 \\
LongPhase & 8 & 2 & 0 & 6 \\
WhatsHap & 33 & 21 & 2 & 10 \\
WhatsHap-dg & 17 & 11 & 0 & 6 \\
longcallD & 15 & 5 & 4 & 6 \\
\multicolumn{5}{@{}l}{\emph{chr6}} \\
pgphase & 9 & 1 & 0 & 8 \\
HiPhase & 36 & 22 & 2 & 12 \\
LongPhase & 14 & 6 & 0 & 8 \\
WhatsHap & 61 & 35 & 8 & 18 \\
WhatsHap-dg & 40 & 22 & 2 & 16 \\
longcallD & 43 & 25 & 7 & 11 \\
\multicolumn{5}{@{}l}{\emph{chr16}} \\
pgphase & 9 & 1 & 0 & 8 \\
HiPhase & 21 & 12 & 0 & 9 \\
LongPhase & 10 & 6 & 0 & 4 \\
WhatsHap & 46 & 33 & 2 & 11 \\
WhatsHap-dg & 21 & 15 & 2 & 4 \\
longcallD & 27 & 17 & 3 & 7 \\
\multicolumn{5}{@{}l}{\emph{chr21}} \\
pgphase & 8 & 1 & 2 & 5 \\
HiPhase & 10 & 6 & 0 & 4 \\
LongPhase & 9 & 4 & 0 & 5 \\
WhatsHap & 18 & 5 & 5 & 8 \\
WhatsHap-dg & 11 & 4 & 3 & 4 \\
longcallD & 13 & 3 & 5 & 5 \\
\bottomrule
\end{longtable}
}

\section{Running time and memory}

{\footnotesize
\setlength{\LTcapwidth}{\linewidth}
\begin{longtable}{@{}lrrr@{}}
\caption{Wall-clock time and peak memory per chromosome and tool.}\label{tab:resources}\\
\toprule
Tool & Threads & Time (s) & Peak memory (GB) \\
\midrule
\endfirsthead
\caption[]{(continued)}\\
\toprule
Tool & Threads & Time (s) & Peak memory (GB) \\
\midrule
\endhead
\multicolumn{4}{@{}l}{\emph{chr12}} \\
pgphase & 6 & 113 & 6.52 \\
HiPhase & 20 & 203 & 0.83 \\
LongPhase & 20 & 32 & 2.55 \\
WhatsHap & 1 & 635 & 1.38 \\
WhatsHap-dg & 1 & 715 & 1.71 \\
longcallD & 20 & 84 & 9.44 \\
\multicolumn{4}{@{}l}{\emph{chr17}} \\
pgphase & 6 & 74 & 5.02 \\
HiPhase & 6 & 158 & 0.33 \\
LongPhase & 6 & 23 & 1.53 \\
WhatsHap & 1 & 422 & 0.90 \\
WhatsHap-dg & 1 & 474 & 1.11 \\
longcallD & 6 & 72 & 13.26 \\
\multicolumn{4}{@{}l}{\emph{chr18}} \\
pgphase & 6 & 67 & 5.31 \\
HiPhase & 20 & 114 & 0.51 \\
LongPhase & 20 & 20 & 1.61 \\
WhatsHap & 1 & 363 & 0.94 \\
WhatsHap-dg & 1 & 405 & 1.18 \\
longcallD & 20 & 66 & 13.10 \\
\multicolumn{4}{@{}l}{\emph{chr19}} \\
pgphase & 6 & 62 & 6.12 \\
HiPhase & 6 & 113 & 0.34 \\
LongPhase & 6 & 19 & 1.23 \\
WhatsHap & 1 & 378 & 0.93 \\
WhatsHap-dg & 1 & 427 & 1.26 \\
longcallD & 6 & 54 & 4.73 \\
\multicolumn{4}{@{}l}{\emph{chr20}} \\
pgphase & 6 & 64 & 6.92 \\
HiPhase & 20 & 106 & 0.48 \\
LongPhase & 20 & 16 & 1.32 \\
WhatsHap & 1 & 301 & 0.82 \\
WhatsHap-dg & 1 & 337 & 1.04 \\
longcallD & 20 & 43 & 11.39 \\
\multicolumn{4}{@{}l}{\emph{chr6}} \\
pgphase & 6 & 147 & 7.60 \\
HiPhase & 6 & 459 & 0.48 \\
LongPhase & 6 & 92 & 3.38 \\
WhatsHap & 1 & 1004 & 1.69 \\
WhatsHap-dg & 1 & 1113 & 2.05 \\
longcallD & 6 & 235 & 8.48 \\
\multicolumn{4}{@{}l}{\emph{chr16}} \\
pgphase & 6 & 89 & 5.87 \\
HiPhase & 6 & 280 & 0.49 \\
LongPhase & 6 & 48 & 1.91 \\
WhatsHap & 1 & 491 & 1.09 \\
WhatsHap-dg & 1 & 548 & 1.45 \\
longcallD & 6 & 80 & 12.40 \\
\multicolumn{4}{@{}l}{\emph{chr21}} \\
pgphase & 6 & 56 & 8.46 \\
HiPhase & 6 & 91 & 0.27 \\
LongPhase & 6 & 16 & 0.98 \\
WhatsHap & 1 & 219 & 0.63 \\
WhatsHap-dg & 1 & 246 & 0.83 \\
longcallD & 6 & 36 & 7.65 \\
\bottomrule
\end{longtable}
}

\section{Scoring settings}

\begin{table}[h]
\caption{Scoring settings per chromosome: the interval counted as centromeric and pericentromeric (CHM13 cenSat v2.1 satellite span on chr6, chr16, chr17, chr19 and chr21, the acrocentric short arm included on chr21; rounded spans on chr12, chr18 and chr20, set earlier in development), and the regions not scored because the truth contradicts itself there.}\label{tab:scoring}
\centering\footnotesize
\setlength{\tabcolsep}{4pt}
\begin{tabular}{@{}lll@{}}
\toprule
Chromosome & Centromeric interval (Mb) & Not scored \\
\midrule
chr12 & 34.70--37.20 & none \\
chr17 & 21.82--28.14 & none \\
chr18 & 15.40--20.90 & none \\
chr19 & 23.83--30.40 & none \\
chr20 & 26.00--32.00 & chr20:66.10--66.30~Mb \\
chr6 & 57.22--62.67 & none \\
chr16 & 33.95--52.26 & none \\
chr21 & 0.00--12.42 & none \\
MHC (chr6:28.5--34.0~Mb) & 57.22--62.67 & none \\
\bottomrule
\end{tabular}
\end{table}

\section{Parameters}

{\footnotesize
\setlength{\LTcapwidth}{\linewidth}
\begin{longtable}{@{}>{\raggedright\arraybackslash}p{5.0cm}rrl@{}}
\caption{Every threshold of pgphase, with its value as compiled (read from the source when this table is generated), the Methods section that uses it, and its name in the source.}\label{tab:params}\\
\toprule
Parameter & Value & Section & Name \\
\midrule
\endfirsthead
\caption[]{(continued)}\\
\toprule
Parameter & Value & Section & Name \\
\midrule
\endhead
\multicolumn{4}{@{}l}{\emph{Graph evidence and catalog sites}} \\
Minimum mapping quality of graph alignments without a BAM & 5 & 2.2 & \texttt{kDefaultGraphMinMapq} \\
Minimum mapping quality of graph and linear alignments with a BAM & 1 & 2.2 & \texttt{kUnionAlignmentMinMapq} \\
Chunk length without a BAM (bp) & 500{,}000 & 2.1 & \texttt{kDefaultChunkSize} \\
Chunk length with a BAM (bp) & 1{,}000,000 & 2.1 & \texttt{kDefaultGraphBamChunkSize} \\
Minimum reads at a site & 5 & 2.2 & \texttt{kDefaultMinDepth} \\
Minimum reads carrying the alternative allele & 2 & 2.2 & \texttt{kDefaultMinAltDepth} \\
Minimum alternative allele fraction of a heterozygous site & 0.2 & 2.2 & \texttt{kDefaultMinAf} \\
Maximum alternative allele fraction of a heterozygous site & 0.8 & 2.2 & \texttt{kDefaultMaxAf} \\
Minimum base quality of an aligned base & 10 & 2.3 & \texttt{kDefaultMinBaseq} \\
Family-wise error of the test for an absent reference base at catalog SNPs & 0.01 & 2.3 & \texttt{kFamilywiseError} \\
Deletion share of reads that keeps a catalog SNP from being reclassified & 0.2 & 2.3 & \texttt{kMinDeletionFraction} \\
\multicolumn{4}{@{}l}{\emph{The sample's own heterozygous variants}} \\
Mapping quality of a confident read & 30 & 2.4 & \texttt{kDefaultMinMapq} \\
Confident reads needed on each allele of an MSA-verified variant & 2 & 2.4 & \texttt{kUnionMinConfidentAlleleReads} \\
Largest depth, relative to the chunk median, of an added site & 1.5 & 2.4 & \texttt{kUnionMaxDepthRatio} \\
Longest noisy region resolved by MSA (bp) & 2{,}000 & 2.4 & \texttt{kUnionMsaMaxRegionLength} \\
Longest catalog site whose calls only fill reads without a graph allele (bp) & 2{,}000 & 2.4 & \texttt{kUnionMaxComplementRowSpan} \\
Minimum mapping quality of reads only the linear alignment holds & 20 & 2.4 & \texttt{kAddedReadMinMapq} \\
Slope \(\lambda\) of the per-read error \(1/(1+e^{\lambda m})\) of a realigned call & 3 & 2.4 & \texttt{kMarginErrorLambda} \\
Error of an MSA call that realignment cannot decide & 0.25 & 2.4 & \texttt{kMsaUndecidedError} \\
Error of an MSA call that realignment contradicts & 0.35 & 2.4 & \texttt{kMsaContradictedError} \\
Error of a graph call at a catalog SNP whose base the alignment does not give & 0.25 & 2.4 & \texttt{kGraphOnlyCallError} \\
\multicolumn{4}{@{}l}{\emph{Expectation-maximization}} \\
Starting error of clean sites & 0.05 & 2.5 & \texttt{kEmCleanStartError} \\
Starting error of other sites & 0.2 & 2.5 & \texttt{kEmOtherStartError} \\
Smallest site error & 0.01 & 2.5 & \texttt{kEmMinError} \\
Largest site and call error & 0.45 & 2.5 & \texttt{kEmMaxError} \\
Pseudo-reads of the class prior & 5 & 2.5 & \texttt{kEmClassPriorReads} \\
Teaching weight (reads) for a site to enter its class pool & 10 & 2.5 & \texttt{kEmPoolMinWeight} \\
Sites needed for the pool median (else the weighted mean) & 5 & 2.5 & \texttt{kEmPoolMinSites} \\
Minimum mapping quality of a read that updates the model & 20 & 2.5 & \texttt{kEmLearnMinMapq} \\
Minimum share of the rarer allele at a site & 0.25 & 2.5 & \texttt{kEmMinMinorAlleleFraction} \\
Iterations before switch moves & 15 & 2.5 & \texttt{kEmIterations} \\
Most rounds of switch moves & 30 & 2.5 & \texttt{kEmSwitchRounds} \\
Learned error at which a site is unphased and labels no read & 0.3 & 2.5 & \texttt{kEmUnreliableError} \\
Posterior for a read label & 0.8 & 2.5 & \texttt{kEmReadLabelPosterior} \\
\multicolumn{4}{@{}l}{\emph{Joins}} \\
Join margin below which a block is cut (log units) & 4 & 2.6 & \texttt{kEmBlockCutLogLikelihood} \\
Reads with evidence on both sides of a join & 2 & 2.6 & \texttt{kEmMinJoinReads} \\
Margin needed by a join through weak sites (log units) & 10 & 2.6 & \texttt{kEmWeakBridgeCutLogLikelihood} \\
Reads at reliable, non-weak sites for a strong margin & 4 & 2.6 & \texttt{kEmMinStrongJoinReads} \\
Callers, relative to the chunk median, above which a site joins nothing & 2 & 2.6 & \texttt{kEmJoinMaxDepthRatio} \\
Callers, relative to the median, of a duplicated-haplotype site & 1.3 & 2.6 & \texttt{kEmDupDepthRatio} \\
Rarer-allele share below which such a site joins nothing & 0.4 & 2.6 & \texttt{kEmDupMinorFraction} \\
Mean mapping quality of a site's callers below which it joins nothing & 40 & 2.6 & \texttt{kEmJoinMinMeanMapq} \\
Distance within which a site can repeat another (bp) & 30 & 2.6 & \texttt{kEmEchoMaxGap} \\
Share of reads two such sites must have in common & 0.8 & 2.6 & \texttt{kEmEchoMinShare} \\
Share of shared reads that must split the same way & 0.95 & 2.6 & \texttt{kEmEchoMinConcordance} \\
Minimum shared reads & 4 & 2.6 & \texttt{kEmEchoMinReads} \\
\multicolumn{4}{@{}l}{\emph{Labels for the remaining reads, and output}} \\
Noisy calls merged into one local window within (bp) & 100 & 2.7 & \texttt{kLocusMergeGap} \\
Padding of a local window (bp) & 40 & 2.7 & \texttt{kLocusPad} \\
Longest tandem repeat period a window extends through (bp) & 50 & 2.7 & \texttt{kLocusMaxPeriod} \\
Longest repeat scan (bp) & 500 & 2.7 & \texttt{kLocusMaxRepeatScan} \\
Longest local window (bp) & 2{,}000 & 2.7 & \texttt{kLocusMaxLength} \\
Most reads per local window & 400 & 2.7 & \texttt{kLocusMaxReads} \\
Minimum mapping quality of a window's reads & 5 & 2.7 & \texttt{kLocusMinMapq} \\
Agreement of a read's phased sites for it to seed a haplotype & 0.9 & 2.7 & \texttt{kSeedLabelMinAgreement} \\
Seeded reads per haplotype medoid & 10 & 2.7 & \texttt{kLocusSeedSample} \\
Rounds of the unseeded split & 3 & 2.7 & \texttt{kLocusSplitRounds} \\
Edits by which a read must be closer to one haplotype or allele & 2 & 2.7 & \texttt{kLocusMinMargin} \\
Reads needed on the rarer side of a window & 3 & 2.7 & \texttt{kLocusMinSide} \\
Share of reads needed on the rarer side & 0.2 & 2.7 & \texttt{kLocusMinMinorFraction} \\
Largest per-read error of a window call & 0.45 & 2.7 & \texttt{kLocusMaxObsError} \\
Flank of an allele window beyond its repeat (bp) & 20 & 2.7 & \texttt{kAlleleWindowFlank} \\
Most reads per allele window & 400 & 2.7 & \texttt{kAlleleWindowMaxReads} \\
Largest reference share of an indel in a two-allele window & 0.2 & 2.7 & \texttt{kPairMaxRefFraction} \\
Shortest tandem repeat tract of a length window (bp) & 6 & 2.7 & \texttt{kLengthMinTract} \\
Spanning reads needed to genotype a length window & 8 & 2.7 & \texttt{kLengthMinReads} \\
Two lengths over one, for a heterozygous length window (log units) & 6 & 2.7 & \texttt{kLengthMinHetLLR} \\
A read's least likelihood ratio between the two lengths & 0.5 & 2.7 & \texttt{kLengthMinReadLLR} \\
Reads by which one orientation must lead to join chunks & 3 & 2.7 & \texttt{kStitchMinNetReads} \\
Share of reads that orientation needs & 0.9 & 2.7 & \texttt{kStitchMinShare} \\
P-value of the association that orients an excluded site & 0.01 & 2.7 & \texttt{kRescueSiteOrientationPValue} \\
Upper 95\% bound on discordance for a site oriented by one locus & 0.15 & 2.7 & \texttt{kRescueSingletonMaxDiscordance} \\
Posterior for a deferred label & 0.7 & 2.7 & \texttt{kEmDeferredLabelPosterior} \\
Largest learned error of an unverified call that labels reads & 0.1 & 2.7 & \texttt{kEmUnverifiedReportError} \\
Posterior for a label from unverified calls & 0.8 & 2.7 & \texttt{kEmUnverifiedLabelPosterior} \\
Largest disagreement of an island site with labelled reads & 0.3 & 2.7 & \texttt{kIslandOutsideMaxError} \\
Unlabelled reads an island site must call & 2 & 2.7 & \texttt{kIslandMinReads} \\
Span of one locus for an unphased phase set (bp) & 10 & 2.7 & \texttt{kOneLocusSpan} \\
Learned error from which a SNP is written unphased & 0.1 & 2.7 & \texttt{kEmVcfMaxSnpError} \\
Learned error from which an indel is written unphased & 0.25 & 2.7 & \texttt{kEmVcfMaxIndelError} \\
\bottomrule
\end{longtable}
}

\section{Commands}

Each tool was run on the reads of one chromosome, CHM13 contig names without the graph's \texttt{CHM13\#0\#} prefix. \texttt{REF} is the chromosome's CHM13 sequence, \texttt{BAM} the surjected reads, \texttt{VCF} the DeepVariant calls, \texttt{T} the threads.

\begin{enumerate}
  \item DeepVariant 1.10.0: \texttt{run\_deepvariant -{}-model\_type=PACBIO -{}-ref=REF -{}-reads=BAM -{}-regions=CHROM -{}-output\_vcf=VCF -{}-num\_shards=T}
  \item HiPhase 1.6.0: \texttt{hiphase -{}-reference REF -{}-bam BAM -{}-vcf VCF -{}-threads T -{}-output-bam phased.bam -{}-output-vcf phased.vcf.gz}
  \item LongPhase 2.0.2: \texttt{longphase phase -s VCF -b BAM -r REF -t T -{}-pb -o phased}, then \texttt{longphase haplotag -s phased.vcf.gz -b BAM -r REF -t T -o phased}
  \item WhatsHap 2.8: \texttt{whatshap phase -{}-reference REF -{}-ignore-read-groups -{}-indels [-{}-distrust-genotypes] -o phased.vcf.gz VCF BAM}, then \texttt{whatshap haplotag -{}-reference REF -o phased.bam phased.vcf.gz BAM}
  \item longcallD v0.0.11: \texttt{longcallD call -{}-hifi -o phased.vcf -b phased.bam -t T REF BAM CHROM}
  \item pgphase: \texttt{pgphase collect-graph-variation -{}-ref REF -{}-bam BAM -{}-sites CATALOG -{}-gaf GAF -r CHROM -{}-threads T -o candidates.tsv -{}-phased-vcf-out phased.vcf -{}-phased-bam-out phased.bam}
  \item Catalog and GAF index: \texttt{pgphase build-snarl-catalog} on the GBZ graph; \texttt{vg convert -G} of the reads' GAM, then \texttt{pgphase index-gaf -{}-graph GBZ -{}-reference-sample CHM13}
  \item Truth labels: \texttt{minimap2 -ax lr:hqae -Y -{}-cs -{}-eqx -e 100 -{}-secondary=no} of the reads to the maternal and, separately, the paternal haplotype of HG002 v1.1; then \texttt{hiphap -p -1 mat -2 pat -A 2 MATERNAL.sam PATERNAL.sam}, which compares the two alignments of each read and assigns it to one haplotype (HipHap aligns nothing itself)
\end{enumerate}

\end{document}
